% ============================================================
% KATA PENGANTAR
% ============================================================
\newpage
\vspace{3cm}
\begin{center}
{\fontsize{14}{16.8}\selectfont\textbf{KATA PENGANTAR}}\\[1em]
\end{center}
\vspace{1em}

\setlength{\parskip}{0.5\baselineskip}

Puji dan syukur penulis panjatkan ke hadirat Tuhan Yang Maha Esa, karena atas berkat, rahmat, dan karunia-Nya, penulis dapat menyelesaikan tesis yang berjudul Restorasi Dokumen Terdegradasi Menggunakan Generative Adversarial Network dengan Diskriminator Dual-Modal dan Optimasi Loss Function Berorientasi HTR.

Penulisan tesis ini dilakukan dalam rangka memenuhi salah satu syarat untuk mencapai gelar Magister pada Program Smart-X Sekolah Teknik Elektro dan Informatika (STIE) Institut Teknologi Bandung.

Penelitian ini dilatarbelakangi oleh kepedulian penulis terhadap tantangan preservasi digital warisan sejarah bangsa, khususnya koleksi dokumen paleografi abad ke-16 hingga ke-18 yang tersimpan di Arsip Nasional Republik Indonesia (ANRI). Penulis menyadari bahwa restorasi dokumen tidak hanya tentang perbaikan visual semata, tetapi juga bagaimana menjembatani kualitas visual tersebut dengan keterbacaan mesin, untuk mendukung digitalisasi arsip dan diperolehnya fakta sejarah baru dari koleksi arsip yang berhasil di rekognisi. Melalui pengembangan kerangka kerja GAN dengan strategi Frozen Recognizer, penelitian ini diharapkan dapat memberikan sumbangsih teknis bagi kemajuan teknologi restorasi dokumen dan pengenalan teks tulisan tangan (HTR) pada dokumen historis

Penulis menyadari bahwa penyelesaian tesis ini tidak terlepas dari bantuan, bimbingan, dan dukungan dari berbagai pihak. Oleh karena itu, pada kesempatan ini, penulis ingin menyampaikan rasa terima kasih dan penghargaan yang setinggi-tingginya kepada:

\begin{enumerate}[leftmargin=1.5em, labelsep=0.5em, itemindent=0pt, noitemsep, topsep=0pt]
    \item Bapak Ir. I Gusti Bagus Baskara Nugraha, ST, MT, Ph.D , selaku pembimbing utama, yang telah meluangkan waktu, tenaga, dan pikiran untuk memberikan arahan strategis, diskusi mendalam mengenai topik penelitian, serta motivasi tanpa henti hingga penelitian ini dapat diselesaikan.

    \item Kementerian Komunikasi dan Digital Republik Indonesia (KOMDIGI) yang telah memberikan beasiswa pasca sarjana, program magister multidisiplin Smart System (Smart-X). 
        
    \item Arsip Nasional Republik Indonesia (ANRI), yang telah memfasilitasi akses terhadap dataset dokumen historis dan seluruh sumber daya komputasi untuk mendukung penelitian ini.
    
    \item Segenap Dosen dan Staf Pengajar di Program Pascasarjana STIE ITB yang telah memberikan bekal ilmu pengetahuan yang sangat bermanfaat selama masa perkuliahan.
    
    \item Seluruh Keluarga, atas doa yang tidak pernah putus, kasih sayang, serta dukungan moral dan material yang menjadi sumber kekuatan penulis dalam melewati berbagai tantangan selama penelitian.
    
    \item Rekan-rekan seperjuangan program Smart-X, terima kasih atas diskusi produktif, dan berbagi pengetahuan, serta kebersamaan yang menyenangkan.
\end{enumerate}

Penulis menyadari bahwa tesis ini masih jauh dari kesempurnaan, mengingat kompleksitas tantangan dalam restorasi dokumen kuno. Oleh karena itu, penulis sangat mengharapkan kritik dan saran yang membangun demi penyempurnaan penelitian ini di masa mendatang.

Akhir kata, penulis berharap semoga tesis ini dapat memberikan manfaat bagi pengembangan ilmu pengetahuan, khususnya di bidang kecerdasan artifisial dan preservasi digital, serta dapat diaplikasikan untuk pelestarian sejarah bangsa.

\vspace{2cm}

\begin{flushright}
    Bandung,   Desember 2025 \\
    \vspace{2cm}
    \textbf{Jatniko Nur Mutaqin}
\end{flushright}

\setlength{\parskip}{0pt}

% --- SELESAI KATA PENGANTAR ---